% Fichier généré automatiquement par tools/generate_corrections.py.
% Source : SLCI/SLCI-03-SchemaBlocs/78_RobotDaVinci/78_RobotDaVinci.tex
% Statut : complet (3/3 question(s) renseignée(s), 0 reprise(s)).
\begin{corrigebox}[Corrigé extrait de la banque]
\CorrigeQuestion{1}
\begin{center}
\includegraphics[width=\linewidth]{cor_02}

\includegraphics[width=\linewidth]{cor_03}

\end{center}
 

En utilisant l'équation électrique du MCC, on a 
$U_1(p)=\left(L p   + R\right)I_1(p)+E_1(p)$. En utilisant le schéma-blocs :  $I_1(p)=\left(U_1(p) - E(p)\right) D(p)$. On a donc 
$I_1(p)=\dfrac{U_1(p) - E(p)}{R+Lp}$ et $D(p) = \dfrac{1}{R+Lp}$.

En utilisant la première relation de comportement du MCC, on a $E_1(p)$ en sortie du bloc $k_e$ et $p\Delta_1(p)$ en entrée; donc $H(p)=\dfrac{1}{p}$.

En utilisant la seconde relation, on a $F(p)=k_t$.

En utilisant l'équation de mouvement de l'axe 1, on a :
$\Delta C_1(p)=J p ^2  \Delta \theta_1(p) - k_1 \dfrac{r_9'}{r_0}h_2 \Delta F_x(p)$.

D'après le schéma-blocs, on a $\Delta \theta_1(p) = \left(\Delta C_1(p)+\Delta F_x(p) E(p)\right) G(p)H(p)$.

En réageançant l'équation, on a 
$J p ^2  \Delta \theta_1(p) = \Delta C_1(p) +  k_1 \dfrac{r_9'}{r_0}h_2 \Delta F_x(p) $
$ \Leftrightarrow   \Delta \theta_1(p) = \left(\Delta C_1(p) +  k_1 \dfrac{r_9'}{r_0}h_2 \Delta F_x(p)\right) \dfrac{1}{J p ^2} $.

On a donc $E(p)= k_1 \dfrac{r_9'}{r_0}h_2 $. 

De plus $G(p)H(p)=\dfrac{1}{Jp^2}$ et $H(p)=\dfrac{1}{p}$; donc $G(p)=\dfrac{1}{Jp}$.

En utilisant l'équation électrique du MCC, on a 
$U_1(p)=\left(L p   + R\right)I_1(p)+E_1(p)$. En utilisant le schéma-blocs :  $I_1(p)=\left(U_1(p) - E(p)\right) D(p)$. On a donc 
$I_1(p)=\dfrac{U_1(p) - E(p)}{R+Lp}$ et $D(p) = \dfrac{1}{R+Lp}$.

En utilisant l'équation du PID, on a 
$U_1(p)=\left( \Delta \theta_{c1}(p)-\Delta \theta_1(p)\right) \left(\sigma_1 + \dfrac{\sigma_2}{p}\right)- \sigma_3p \Delta \theta_1(p)+\sigma_4\Delta \theta_{c1}(p)$
soit $U_1(p)=\left( \Delta \theta_{c1}(p) \left(\sigma_1 + \dfrac{\sigma_2}{p}\right) -\Delta \theta_1(p) \left(\sigma_1 + \dfrac{\sigma_2}{p}\right)\right) - \sigma_3p \Delta \theta_1(p)+\sigma_4\Delta \theta_{c1}(p)$.

En utilisant le schéma-blocs, on a 
$U_1(p)=\Delta_{c1}(p) A(p) + \left(\Delta_{c1}(p) -\Delta\theta_{1}(p)\right) B(p) - \Delta\theta_{1}(p)C(p)$
$=\Delta_{c1}(p) \left( A(p) + B(p)\right) -\Delta\theta_{1}(p)  \left(B(p)+C(p)\right)$.

Par suite, 
$U_1(p)= \Delta \theta_{c1}(p) \left(\sigma_1 + \dfrac{\sigma_2}{p} +\sigma_4\right) -\Delta \theta_1(p) \left(\sigma_1 + \dfrac{\sigma_2}{p} + \sigma_3p\right)$.

On aura donc $B(p)=\sigma_1 + \dfrac{\sigma_2}{p}$, $C(p)=\sigma_3 p$ et  $A(p)=\sigma_4$.

En utilisant l'équation électrique du MCC, on a 
$U_1(p)=\left(L p   + R\right)I_1(p)+E_1(p)$. En utilisant le schéma-blocs :  $I_1(p)=\left(U_1(p) - E(p)\right) D(p)$. On a donc 
$I_1(p)=\dfrac{U_1(p) - E(p)}{R+Lp}$ et $D(p) = \dfrac{1}{R+Lp}$.

En utilisant la première relation de comportement du MCC, on a $E_1(p)$ en sortie du bloc $k_e$ et $p\Delta_1(p)$ en entrée; donc $H(p)=\dfrac{1}{p}$.

En utilisant la seconde relation, on a $F(p)=k_t$.

En utilisant l'équation de mouvement de l'axe 1, on a :
$\Delta C_1(p)=J p ^2  \Delta \theta_1(p) - k_1 \dfrac{r_9'}{r_0}h_2 \Delta F_x(p)$.

D'après le schéma-blocs, on a $\Delta \theta_1(p) = \left(\Delta C_1(p)+\Delta F_x(p) E(p)\right) G(p)H(p)$.

En réageançant l'équation, on a 
$J p ^2  \Delta \theta_1(p) = \Delta C_1(p) +  k_1 \dfrac{r_9'}{r_0}h_2 \Delta F_x(p) $
$ \Leftrightarrow   \Delta \theta_1(p) = \left(\Delta C_1(p) +  k_1 \dfrac{r_9'}{r_0}h_2 \Delta F_x(p)\right) \dfrac{1}{J p ^2} $.

On a donc $E(p)= k_1 \dfrac{r_9'}{r_0}h_2 $. 

De plus $G(p)H(p)=\dfrac{1}{Jp^2}$ et $H(p)=\dfrac{1}{p}$; donc $G(p)=\dfrac{1}{Jp}$.

En utilisant l'équation électrique du MCC, on a 
$U_1(p)=\left(L p   + R\right)I_1(p)+E_1(p)$. En utilisant le schéma-blocs :  $I_1(p)=\left(U_1(p) - E(p)\right) D(p)$. On a donc 
$I_1(p)=\dfrac{U_1(p) - E(p)}{R+Lp}$ et $D(p) = \dfrac{1}{R+Lp}$.

En utilisant l'équation du PID, on a 
$U_1(p)=\left( \Delta \theta_{c1}(p)-\Delta \theta_1(p)\right) \left(\sigma_1 + \dfrac{\sigma_2}{p}\right)- \sigma_3p \Delta \theta_1(p)+\sigma_4\Delta \theta_{c1}(p)$
soit $U_1(p)=\left( \Delta \theta_{c1}(p) \left(\sigma_1 + \dfrac{\sigma_2}{p}\right) -\Delta \theta_1(p) \left(\sigma_1 + \dfrac{\sigma_2}{p}\right)\right) - \sigma_3p \Delta \theta_1(p)+\sigma_4\Delta \theta_{c1}(p)$.

En utilisant le schéma-blocs, on a 
$U_1(p)=\Delta_{c1}(p) A(p) + \left(\Delta_{c1}(p) -\Delta\theta_{1}(p)\right) B(p) - \Delta\theta_{1}(p)C(p)$
$=\Delta_{c1}(p) \left( A(p) + B(p)\right) -\Delta\theta_{1}(p)  \left(B(p)+C(p)\right)$.

Par suite, 
$U_1(p)= \Delta \theta_{c1}(p) \left(\sigma_1 + \dfrac{\sigma_2}{p} +\sigma_4\right) -\Delta \theta_1(p) \left(\sigma_1 + \dfrac{\sigma_2}{p} + \sigma_3p\right)$.

On aura donc $B(p)=\sigma_1 + \dfrac{\sigma_2}{p}$, $C(p)=\sigma_3 p$ et  $A(p)=\sigma_4$.
\CorrigeQuestion{2}
On a $H_{\text{processus}}(p)= \dfrac{D(p)F(p)G(p)}{1+D(p)F(p)G(p) k_e} H(p)$
soit $H_{\text{processus}}(p)= \dfrac{\dfrac{1}{R+Lp}k_t\dfrac{1}{Jp}}{1+\dfrac{1}{R+Lp}k_t\dfrac{1}{Jp} k_e} \dfrac{1}{p}$.
Avec $L=0$, 
$H_{\text{processus}}(p) =\dfrac{k_t}{RJp+k_t k_e} \dfrac{1}{p}=\dfrac{\dfrac{1}{k_e}}{\dfrac{RJ}{k_t k_e}p+1} \dfrac{1}{p}$ soit $K = \dfrac{1}{k_e}$ et $\tau = \dfrac{RJ}{k_t k_e}$.
\CorrigeQuestion{3}
On a vu que $U_1(p)= \Delta \theta_{c1}(p) \left(\sigma_1 + \dfrac{\sigma_2}{p} +\sigma_4\right) -\Delta \theta_1(p) \left(\sigma_1 + \dfrac{\sigma_2}{p} + \sigma_3p\right)$ et 
que $\dfrac{\Delta \theta_1(p)}{U_1(p)}=\dfrac{K}{p\left(1+\tau p \right)}$.

On a donc 
$\Delta \theta_1(p) \dfrac{p\left(1+\tau p \right)}{K}= \Delta \theta_{c1}(p) \left(\sigma_1 + \dfrac{\sigma_2}{p} +\sigma_4\right) -\Delta \theta_1(p) \left(\sigma_1 + \dfrac{\sigma_2}{p} + \sigma_3p\right)$

$\Leftrightarrow \Delta \theta_1(p) \left(\dfrac{p\left(1+\tau p \right)}{K}  + \sigma_1 + \dfrac{\sigma_2}{p} + \sigma_3p\right)= \Delta \theta_{c1}(p) \left(\sigma_1 + \dfrac{\sigma_2}{p} +\sigma_4\right) $ et 

$B_F(p) = \dfrac{\sigma_1 + \dfrac{\sigma_2}{p} +\sigma_4}{\dfrac{p\left(1+\tau p \right)}{K}  + \sigma_1 + \dfrac{\sigma_2}{p} + \sigma_3p}$ 
$= \dfrac{\sigma_1 p + \sigma_2 + \sigma_4 p}{\dfrac{p^2\left(1+\tau p \right)}{K}  + \sigma_1p + \sigma_2 + \sigma_3p^2}$
$= K \dfrac{\sigma_1 p + \sigma_2 + \sigma_4 p}{p^2\left(1+\tau p \right)  + \sigma_1 K p + \sigma_2 K  + \sigma_3 Kp^2}$
$= K \dfrac{\left(\sigma_1+ \sigma_4 \right) p + \sigma_2 }{ \tau p^3 + p^2\left(1+\sigma_3 \right)  + \sigma_1 K p + \sigma_2 K  }$.

\begin{enumerate}
\item 
$A(p)=\sigma_4$,
$B(p)=\sigma_1 + \dfrac{\sigma_2}{p}$, 
$C(p)=\sigma_3 p$, 
$D(p) = \dfrac{1}{R+Lp}$, 
$E(p)= k_1 \dfrac{r_9'}{r_0}h_2 $, 
$F(p)=k_t$, 
$G(p)=\dfrac{1}{Jp}$, 
$H(p)=\dfrac{1}{p}$.
\item $K = \dfrac{1}{k_e}$ et $\tau = \dfrac{RJ}{k_t k_e}$.
\item $B_F(p) = K \dfrac{\left(\sigma_1+ \sigma_4 \right) p + \sigma_2 }{ \tau p^3 + p^2\left(1+\sigma_3 \right)  + \sigma_1 K p + \sigma_2 K  }$.
\end{enumerate}
\end{corrigebox}
